\banner{Design project 11: a furrow system for 250-m runs}\label{sec:proj11}
\textbf{Brief.} A furrow system on a clay loam of slope 0.002 is to deliver a net depth of 60 mm in furrows 250 m long and 0.75 m apart. The required opportunity time is 62 min. The designer tries four inflows from 0.9 to 2.0 L\,s$^{-1}$ and compares the advance time, the cut-off time and the efficiency. The advance is solved from a volume balance with surface storage $\sigma_y A_0$ and subsurface storage $\sigma_z k t^a$.

\begin{table}[!htbp]\centering\small\caption{Furrow design alternatives for $L=250$ m, $S_0=0.002$, spacing 0.75 m.}
\begin{tabular}{rrrrr}\toprule
\hd{$q_0$ (L/s)} & \hd{Advance (min)} & \hd{Cut-off (min)} & \hd{Applied depth (mm)} & \hd{$E_a$ (\%)}\\\midrule
0.9 & 691 & 408 & 117 & 51\\
1.2 & 334 & 229 & 88 & 68\\
1.6 & 172 & 148 & 76 & 79\\
2.0 & 107 & 116 & 74 & 81\\\bottomrule\end{tabular}\end{table}

\textbf{Reading the table.} A larger inflow advances faster, so the head of the field is not over-watered while the tail is wet: efficiency rises. The erosive limit of 3.1 L\,s$^{-1}$ caps the choice; the designer takes the largest stable flow that stays inside the limit, and chooses cut-off from the opportunity-time requirement. The tail-water can be reused or cut back.
\fig{g11_adv}{0.7}{Advance of the water front along a 250-m furrow for four inflows, from the volume balance.}
